\element{parfait}{
  \begin{questionmult}{parfait}\bareme{haut=2,p=0}
    Dans un écoulement parfait, on néglige
    % \begin{multicols}{2}
      \begin{reponses}
        \bonne{la diffusion thermique}
        \bonne{la viscosité}
        \mauvaise{le déplacement des particules de fluide}
        \mauvaise{le nombre de Reynolds devant $1$}
      \end{reponses}
    % \end{multicols}
  \end{questionmult}
}
\setgroupmode{parfait}{cyclic}

\element{Bernoulli}{
  \begin{questionmult}{hypBernoulli}\bareme{haut=2,p=0}
    Pour que la relation de Bernoulli soit applicable, il faut
    % \begin{multicols}{2}
      \begin{reponses}
        \bonne{un écoulement incompressible}
        \bonne{un écoulement stationnaire}
        \bonne{un écoulement homogène}
        \bonne{un écoulement parfait}
        \bonne{deux points situés sur une même ligne de courant}
        \mauvaise{un écoulement laminaire}
      \end{reponses}
    % \end{multicols}
  \end{questionmult}
}
\setgroupmode{Bernoulli}{cyclic}

\element{PPI}{
  \begin{question}{PPI1}\bareme{1}
    Le premier principe industriel s'écrit
    % \begin{multicols}{2}
      \begin{reponses}
        \bonne{$h_s+e_{c,s}+e_{p,s} - (h_e+e_{c,e}+e_{p,e}) = w_u + q$}
        \mauvaise{$h_s+e_{c,s}+e_{p,s} - (h_e+e_{c,e}+e_{p,e}) = P_u + P_{th}$}
        \mauvaise{$h(t_2)+e_c(t_2)+e_p(t_2) - (h(t_1)+e_c(t_1)+e_p(t_1))  = P_u + P_{th}$}
        \mauvaise{$h(t_2)+e_c(t_2)+e_p(t_2) - (h(t_1)+e_c(t_1)+e_p(t_1))  = w_u + q$}
      \end{reponses}
    % \end{multicols}
  \end{question}
}
\element{PPI}{
  \begin{question}{PPI2}\bareme{1}
    Le premier principe industriel s'écrit
    % \begin{multicols}{2}
      \begin{reponses}
        \mauvaise{$D_m \left( h_s+e_{c,s}+e_{p,s} - (h_e+e_{c,e}+e_{p,e}) \right) = w_u + q$}
        \bonne{$D_m \left(h_s+e_{c,s}+e_{p,s} - (h_e+e_{c,e}+e_{p,e}) \right) = P_u + P_{th}$}
        \mauvaise{$D_m \left(h(t_2)+e_c(t_2)+e_p(t_2) - (h(t_1)+e_c(t_1)+e_p(t_1)) \right)  = P_u + P_{th}$}
        \mauvaise{$D_m \left(h(t_2)+e_c(t_2)+e_p(t_2) - (h(t_1)+e_c(t_1)+e_p(t_1)) \right)  = w_u + q$}
      \end{reponses}
    % \end{multicols}
  \end{question}
}
\setgroupmode{PPI}{cyclic}

\element{théorèmes}{
  \begin{questionmult}{th}\bareme{haut=2,p=0}
    Parmi les relations suivantes, lesquelles ne s'appliquent qu'à des systèmes ouverts ?
    % \begin{multicols}{2}
      \begin{reponses}
        \mauvaise{théorème de l'énergie mécanique}
        \mauvaise{théorème de la puissance cinétique}
        \mauvaise{théorème de la résultante cinétique}
        \mauvaise{théorème du moment cinétique}
        \mauvaise{premier principe de la thermodynamique}
        \mauvaise{second principe de la thermodynamique}
      \end{reponses}
    % \end{multicols}
  \end{questionmult}
}
\setgroupmode{théorèmes}{cyclic}